\documentclass[11pt,a4paper]{article}
\usepackage[utf8]{inputenc}
\usepackage[T1]{fontenc}
\usepackage[a4paper,margin=2.5cm]{geometry}
\usepackage{microtype}
\usepackage{hyperref}
\usepackage{xurl}
\usepackage{xcolor}
\sloppy
\hypersetup{colorlinks=true, linkcolor=blue, urlcolor=blue, citecolor=blue}

\setlength{\parskip}{0.4em}
\setlength{\parindent}{0pt}

\begin{document}

\textbf{Cover letter to the Editors of \emph{Quantum}}

\vspace{0.4cm}

Hikaru Wakaura, Taiki Tanimae \\
QIRI (Quantum Integrated Research Institute Inc.)\\
1--16--3, Akasaka, Minato-ku, Tokyo 107--0061, Japan\\
\href{mailto:h.wakaura@qiri.co.jp}{h.wakaura@qiri.co.jp} \quad ORCID: 0009-0001-2341-7710

\vspace{0.3cm}
April 28, 2026

\vspace{0.4cm}

Dear Editors,

\medskip

We are pleased to submit our manuscript, \textbf{``Catalytic Quantum Error Correction: Theory, Efficient Catalyst Preparation, and Numerical Benchmarks''}, for consideration by \emph{Quantum}.

\medskip

\textbf{Summary of the contribution.}
We turn a recent foundational theorem in the resource theory of coherence (Shiraishi \& Takagi, \emph{Phys.\ Rev.\ Lett.}\ \textbf{132}, 180202, 2024) into a candidate quantum-error-correction primitive that we call \emph{Catalytic Quantum Error Correction} (CQEC): a protocol with explicit mode-inclusion success conditions that recovers a known target state from noisy copies without an error-\emph{magnitude} threshold.
The constructive proof requires $n^{*} \sim d^{4} e^{2\gamma}$ noisy copies---about $7.3 \times 10^{10}$ for a 6-qubit catalyst at moderate dephasing; our three-stage pipeline (CPMG dynamical decoupling, Clifford twirling, recursive swap-test purification) reaches the same $F_\mathrm{cat} \geq 0.99$ endpoint from 32~copies, a $10^{4}$--$10^{9}$-fold reduction depending on dimension.
Distinctively, the manuscript pairs its fast effective model with a \emph{genuinely CPTP, exactly covariant joint-channel implementation} (system--catalyst--ancilla, catalyst obtained by partial trace), validates real recovery under dephasing at small dimension ($0.54 \to 0.77$), and quantifies the model-vs-channel gap openly---including a target-swap control experiment and the finding that shallow circuits consume rather than preserve the catalyst.

\medskip

\textbf{Why \emph{Quantum}.}
\emph{Quantum}'s emphasis on theoretical rigor, comprehensive numerical evidence, and open scientific practice matches every aspect of this manuscript:
\begin{itemize}
\item \emph{Theoretical framework.}
We formulate the CQEC protocol with explicit success/failure conditions based on the mode inclusion criterion $\mathcal{C}(\rho_{0}) \subseteq \mathcal{C}(\rho_\mathrm{noisy})$, give finite-copy fidelity bounds derived from Ref.\ [Shiraishi \& Takagi], and clarify its relationship to stabilizer QEC, covariant codes (Eastin--Knill), and the concurrent Purification QEC of Raghoonanan \& Byrnes (arXiv:2603.11568).
\item \emph{Comprehensive benchmarks.}
Four quantum algorithms (qDRIFT, QKAN, control-free QPE, Regev factoring) at $d = 4$--$64$, three noise channels (dephasing, depolarizing, combined), 200~configurations, a tree-tensor-network cryptographic protocol, a 100-cycle catalyst-reuse durability test, gate-noise resilience, circuit-depth scaling, and entangled-state recovery (Bell, GHZ, W).
\item \emph{Quantitative claims with explicit scope.}
$F > 0.99$ across the 200-point effective-model sweep, with the model's assumptions and its gap to the CPTP channel stated and measured; mode-inclusion classification verified over 10~orders of magnitude in residual coherence; $F_\mathrm{cat} > 0.96$ from $n = 8$ copies (0.99 from $n = 32$) for $d = 4$--$64$ under the stated DD/twirl model.
\item \emph{Open science.}
The complete reference implementation is released as the \texttt{cqec} Python package v0.2.0 (MIT license) with a 52-test \texttt{pytest} suite---including tests that lock in the joint channel's linearity, complete positivity, exact covariance, and catalyst-by-partial-trace properties---and a one-command reproduction script: \href{https://github.com/deeptell-inc/cqec}{\texttt{github.com/deeptell-inc/cqec}}.
\end{itemize}

\medskip

\textbf{Key results highlighted for editors.}
\begin{enumerate}
\item \emph{CPTP joint-channel validation.} A genuinely linear, completely positive, exactly covariant recovery channel ($\|[U,H]\| = 0$, linearity residual $< 2 \times 10^{-16}$, Choi-positive) improves fidelity under dephasing at $d = 4$ from $0.54$ to $0.77$; a target-swap control shows the recovery is a property of the channel, not of injected target information.
\item \emph{Honest gap analysis.} The same channel does not yet improve fidelity under depolarizing noise or at $d = 8$, and its catalyst marginal degrades ($0.48$--$0.84$)---the manuscript reports this openly and identifies deeper joint circuits with catalyst-preservation constraints as the research frontier.
\item \emph{Resource bottleneck reduced.} Catalyst preparation at the matched $F_\mathrm{cat} \geq 0.99$ endpoint drops from $7.3 \times 10^{5}$--$7.3 \times 10^{10}$ copies (distillation) to 32~copies, with postselection costs quantified ($P_\mathrm{succ} \approx 0.7$).
\item \emph{Model-sensitivity disclosure.} The pipeline's dimension independence is shown to depend on the twirl model (uniform vs.\ mode-resolved), a sensitivity the paper derives explicitly rather than hiding.
\item \emph{Entanglement restoration (effective model).} Bell-state concurrence $0.135 \to 0.682$ ($5.0\times$) under a global recovery map, with the target-dependence caveat stated.
\end{enumerate}

\medskip

\textbf{Companion preprint and code (all publicly available).}
The manuscript is archived at \href{https://doi.org/10.48550/arXiv.2603.25774}{arXiv:2603.25774}.
A concurrent and complementary submission ``Blind Catalytic Quantum Error Correction'' (treated independently) removes the target-state-knowledge requirement of CQEC; it is also under consideration by \emph{Quantum} and shares no overlapping figures or theorems with the present manuscript.

\medskip

\textbf{Suggested referees.}
We list institutional email addresses obtained from public faculty pages; please verify before contact.
\begin{itemize}
\item \textbf{Naoto Shiraishi} (Graduate School of Arts and Sciences, The University of Tokyo, 3--8--1 Komaba, Meguro-ku, Tokyo 153--8902, Japan; \href{mailto:shiraishi@phys.c.u-tokyo.ac.jp}{\texttt{shiraishi@phys.c.u-tokyo.ac.jp}}) --- co-author of the foundational PRL on catalytic coherence transformations (\emph{Phys.\ Rev.\ Lett.}\ \textbf{132}, 180202, 2024) that our CQEC protocol operationalizes.
\item \textbf{Ryuji Takagi} (School of Physical and Mathematical Sciences, Nanyang Technological University, 21 Nanyang Link, Singapore 637371; \href{mailto:ryuji.takagi@ntu.edu.sg}{\texttt{ryuji.takagi@ntu.edu.sg}}) --- co-author of the same foundational PRL; ideal evaluator for the resource-theoretic content of Sec.~II.
\item \textbf{Tim Byrnes} (NYU Shanghai, 567 West Yangsi Road, Shanghai 200124, China; \href{mailto:tim.byrnes@nyu.edu}{\texttt{tim.byrnes@nyu.edu}}) --- co-author of the concurrent Purification QEC framework (arXiv:2603.11568) directly compared in our Table~IX.
\item \textbf{Iman Marvian} (Department of Electrical and Computer Engineering, Duke University, 100 Science Drive, Durham, NC 27708, USA; \href{mailto:iman.marvian@duke.edu}{\texttt{iman.marvian@duke.edu}}) --- expert in covariant operations, resource theory of asymmetry, and approximate covariant codes (cited in Sec.~II.A).
\item \textbf{Andrew M.~Childs} (Department of Computer Science and Institute for Quantum Information and Computer Science, University of Maryland, College Park, MD 20742, USA; \href{mailto:amchilds@umd.edu}{\texttt{amchilds@umd.edu}}) --- co-author of the recursive swap-test purification protocol (cited as Ref.~\cite{Childs2025}) that the Stage~3 of our pipeline adapts.
\item \textbf{Aram W.~Harrow} (Center for Theoretical Physics, MIT, 77 Massachusetts Ave, Cambridge, MA 02139, USA; \href{mailto:aram@mit.edu}{\texttt{aram@mit.edu}}) --- expert in entanglement distillation, state purification, and related copy-based protocols underlying the swap-test analysis.
\end{itemize}

\textbf{Suggested non-preferred referees.}
None.

\medskip

\textbf{Statements.}
This work has not been published or submitted elsewhere.
All authors have read and approved the submission.
We declare no competing interests.

\medskip

We thank you in advance for your consideration. We hope you find the manuscript suitable for \emph{Quantum} and look forward to the referees' feedback.

\medskip

Sincerely,

\vspace{0.4cm}
Hikaru Wakaura\\
Taiki Tanimae\\
QIRI Inc., Tokyo

\end{document}
